% Standalone appendix document - compiles on its own:
%   latexmk -pdf -outdir=build 06-agenda-2026-10-28.tex
\documentclass[a4paper,11pt]{article}
\usepackage[danish]{babel}
\usepackage{../../appendix-style}
\usepackage{pgfplotstable}
\pgfplotsset{compat=1.18}

\title{12.6 Dagsorden -- 28-10}
\author{SW4PRJ4 -- Group 3}
\date{28-10-2026}

\begin{document}
\maketitle

\section*{Mødeinformation}

\textbf{Dato:} 28-10-2026 (onsdag)

\textbf{Tid:} 10:00

\textbf{Sted:} Incuba

\section*{Dagsordenspunkter}

\begin{enumerate}
    \item \textbf{Sprint review -- sprint 3 (7/10 -- 27/10)} (ca. 45 min)
    \begin{itemize}
        \item Sprint goal fra 6/10 og 7/10: onboarding af en ny bruger og helt styr på Keycloak-flowet, herunder claims, roller og admin-login i admin-panelet. Blev det nået?
        \item Demo: hver af os gennemgår det, vi selv har lavet i sprinten.
        \begin{itemize}
            \item Login i kunde-appen med Keycloak (PMU-70) og onboarding af en ny kunde med navn og adresse gennem \texttt{GET} og \texttt{PUT /user} (PMU-74).
            \item Realm-roller som claims i backend og et admin-endpoint, der kun svarer en admin (PMU-75).
            \item Admin-login i admin-panelet med Keycloak (PMU-76).
            \item Udgivelse af kunde-appen gennem Expo og EAS Build (PMU-82).
            \item Originale kilder i stedet for slides (PMU-81) og gruppesangen (PMU-85).
        \end{itemize}
        \item \textbf{Klar til godkendelse:} PMU-70, PMU-74, PMU-75, PMU-76, PMU-81 og PMU-82 er merget via PR og står i \emph{In Review}. PMU-74 og PMU-76 har ``vist til sprint review'' i deres definition of done. Efter demoen godkender vi dem og flytter dem til \emph{Done}.
        \item PMU-85 (gruppesangen) er merget, men står stadig i \emph{To Do}. Er den færdig?
        \item PMU-76 er færdig uden deployment: DoD-linjen om det deployede admin-panel er flyttet til det nye kort PMU-86 (deploy admin-panelet).
        \item Ikke færdigt: PMU-77 auth-flow (3 point) og PMU-79 Keycloak-features (2 point) står i \emph{In Progress}. PMU-78 flowdiagram for gem-flow (2 point) og PMU-80 zombie-brugere (\emph{? pt}, ikke afklaret) er ikke startet.
        \item Uden for sprinten: PMU-87 (controllere kalder kun services, 1 point) blev merget 28/10 og har ingen sprint-label.
        \item Velocity for sprint 3: committed 31,5 point, done 24,5 point, spillover 7 point, scope 91,5 point. Done forudsætter, at de syv kort ovenfor godkendes. PMU-80 er ikke talt med, fordi det ikke er estimeret. Tallene står i \texttt{docs/assets/data/sprint-velocity.csv}.
        \item Vi blev færdige med 60 \% af det, vi tog ind i sprint 1, 89 \% i sprint 2 og 78 \% i sprint 3. Figur~\ref{fig:velocity} viser de tre sprints ved siden af hinanden.
    \end{itemize}

    \begin{figure}[H]
        \centering
        \pgfplotstableread[col sep=comma]{../../../assets/data/sprint-velocity.csv}\velocitytable
        \begin{tikzpicture}
            \begin{axis}[
                ybar, bar width=14pt,
                width=0.8\textwidth, height=6cm,
                xlabel={Sprint}, ylabel={Point},
                xmin=0.5, xmax=6.5, xtick={1,...,6}, ymin=0,
                legend pos=north east,
                nodes near coords, nodes near coords style={font=\footnotesize},
            ]
                \addplot[fill=blue!35, draw=blue] table[x=sprint, y=committed]{\velocitytable};
                \addlegendentry{Committed}
                \addplot[fill=red!45, draw=red] table[x=sprint, y=done]{\velocitytable};
                \addlegendentry{Done}
            \end{axis}
        \end{tikzpicture}
        \caption{Velocity pr. sprint. Sprint 3 var tre uger med efterårsferie. Det, der ikke blev færdigt, er de tre rapportkort om flows og Keycloak.}
        \label{fig:velocity}
    \end{figure}
    \item \textbf{Retrospective} (ca. 30 min)
    \begin{itemize}
        \item Opfølgning på sidste retrospective: har vi arbejdet test først (PMU-83)? Er målet og testtilfældene skrevet på kortet, før implementeringen startede?
        \item Næsten alt kode blev merget samme dag (21/10). Hvorfor kom det så sent, og hvordan får vi mindre PR'er ind løbende?
        \item Alle tre rapportkort fra vejledermødet er ikke færdige, mens kodekortene er. Prioriterer vi rapporten for lavt?
        \item Kortene blev ikke flyttet på boardet, efterhånden som de blev færdige. Hvordan holder vi boardet ajour?
        \item Hvad gik godt?
        \item Hvad gik skidt?
        \item Hvad ændrer vi i sprint 4? Hver handling bliver et kort i Kaneo.
    \end{itemize}
    \item \textbf{Sprint planning -- sprint 4 (28/10 -- 10/11)} (ca. 60 min)
    \begin{itemize}
        \item Kapacitet: hvor mange timer har hver af os i de to uger? Gennemsnittet af done i sprint 1 -- 3 er 27,8 point, men sprint 3 var en uge længere.
        \item Spillover: PMU-77, PMU-78, PMU-79 og PMU-80 beholder label \emph{Sprint 3} og får \emph{Sprint 4} oveni. Frist flyttes til 10/11.
        \item PMU-80 (zombie-brugere) skal afklares og estimeres, før den kan regnes med. Spørger vi Michel?
        \item PMU-86 (deploy admin-panelet) mangler point-label, epic-label og undergruppe. Skal den med i sprint 4?
        \item PMU-87: tæller vi den med i sprint 4, eller står den uden for sprintene?
        \item \textbf{Backloggen er tom.} Kortene PMU-24 -- PMU-33, PMU-73 og PMU-83 findes ikke længere i Kaneo. Vi skal oprette nye kort ud fra de must-have-krav, der ikke er påbegyndt:
        \begin{itemize}
            \item FR-02: kunder, abonnementer og ordrer gemmes i databasen.
            \item FR-03: standard-pakketyper for voksne og børn.
            \item FR-04: udløbsdato pr. ordre, så næste ordre kan planlægges.
            \item FR-05 og FR-06: admin-panelet viser de næste ordrer, der skal sendes, og admin kan markere en ordre som sendt.
        \end{itemize}
        \item Rapport: hvilke kapitler skal have et udkast i sprint 4?
        \item Planning poker på de nye kort.
        \item Ready-kort fra backlog til \emph{To Do}: \texttt{Krav: FR-xx} som første linje, definition of done, point-label, epic-label og ejer.
        \item Sprint goal (én sætning, der nævner de krav sprinten leverer).
        \item Ejer og undergruppe på hvert kort.
        \item Point i \emph{To Do} (committed).
    \end{itemize}
    \item \textbf{Næste vejledermøde}
    \item \textbf{Eventuelt}
\end{enumerate}

\end{document}
